% TOP_PROBLEM: {"id": 2735, "title": "Kirby Problem 1.76", "category_id": 7, "category": {"id": 7, "name": "topology", "display_name": "Topology", "description": "Properties preserved under continuous deformations.", "slug": "topology", "order_index": 7, "created_at": "2026-07-31T15:26:25.671Z"}, "status": "open", "problem_number": "KP-1.76", "set_id": 11}
% FIRST_POSTED: 2026-10-01
% ENTRY_KIND: statement_only
% UnsolvedMath ID 2735; source catalogue code KP-1.76
% !TeX program = lualatex
% Definition and sources only; this file is not an LLM solution attempt.
\documentclass[11pt,a4paper]{article}
\usepackage[margin=25mm]{geometry}
\usepackage{fontspec}
\setmainfont{lmroman10-regular.otf}[BoldFont=lmroman10-bold.otf,
 ItalicFont=lmroman10-italic.otf,BoldItalicFont=lmroman10-bolditalic.otf]
\usepackage{amsmath,amssymb,mathtools}
\usepackage{unicode-math}
\setmathfont{latinmodern-math.otf}
\AtBeginDocument{\let\setminus\smallsetminus}
\usepackage{xurl,hyperref}
\hypersetup{colorlinks=true,urlcolor=blue}
\setcounter{secnumdepth}{0}
\setlength{\parindent}{0pt}
\setlength{\parskip}{5pt}
\setlength{\emergencystretch}{3em}
\begin{document}
\section*{Kirby Problem 1.76}\label{2735}
{\small UnsolvedMath ID 2735 · KP-1.76 · Topology · Kirby's Problems in Low-Dimensional Topology}
\subsection{Definitions and mathematical statement}
A knot diagram is a generic plane projection of an embedded circle, with
over/under data at its $n$ crossings. It can be encoded by a finite plane graph
with crossing information. The unknot recognition problem returns yes exactly
when this diagram represents the ambient isotopy class of a round circle.

First, is there a deterministic recognition algorithm whose worst-case running
time is bounded by a polynomial in $n$? The polynomial and its degree must be
independent of the input knot. A finite combinatorial encoding is understood,
so arbitrarily complicated numerical coordinates are not additional input data.

Second, determine the structure of the unknot Reidemeister graph. Its vertices
are diagrams of the unknot modulo planar isotopy, and an edge joins two vertices
if one Reidemeister move relates them. The three moves insert or remove a curl,
insert or remove a canceling pair of crossings, or move a strand past a crossing.
Give useful structural descriptions of this infinite connected graph, including
its graph metric and the organization of vertices by crossing number.
The second question allows intermediate diagrams with more crossings than
either endpoint.
\subsection{Short English statement}
Is unknot recognition possible in polynomial time, and what is the structure of the graph of all unknot diagrams and Reidemeister moves?
\subsection{Sources}
\begin{itemize}
\item[S1] ULAM.AI and the UnsolvedMath Contributors, supplied problems.json, record 2735 (KP-1.76), Kirby's Problems in Low-Dimensional Topology. Catalogue identity and source transcription; CC BY 4.0.
\url{https://huggingface.co/datasets/ulamai/UnsolvedMath}.
\item[S2] K3: A New Problem List in Low-Dimensional Topology, Problem 1.76. Expanded formulation of the supplied catalogue statement.
\url{https://math.berkeley.edu/sites/default/files/surv-295-ruberman-watermarked-author-pdf.pdf}.
\end{itemize}
\end{document}
